\begin{figure}[h]
    \centering
    \includegraphics[width=0.9\linewidth]{graphics/integrated_comp/subcontour_supply.pdf}
    \caption{Electricity supply with hourly temporal matching.
    Increasing the hydrogen export volume requires additional renewable electricity supply, as defined in the additionality criteria.
    Advancing domestic \co mitigation leads to the phase-out of coal power, with Combined-Cycle Gas Turbines (CCGT) serving as a transitional substitute before achieving a fossil-free electricity system.
    }
    \label{fig:supply}
\end{figure}


\begin{figure}[h]
    \centering
    \includegraphics[width=0.9\linewidth]{graphics/integrated_comp/subcontour_capacity.pdf}
    \caption{Electricity capacities with hourly temporal matching.
    Increasing the hydrogen export volume requires additional renewable electricity capacities, as defined in the additionality criteria. Fossil capacities remain limited to existing brownfield assets.    
    }
    \label{fig:el-cap}
\end{figure}


\begin{figure}[h]
    \centering
    \includegraphics[width=0.9\linewidth]{graphics/integrated_comp/subcontour_cf.pdf}
    \caption{Electricity capacity factors with hourly temporal matching.
    Coal utilization decreases as domestic \co mitigation progresses, with CCGT acting as a substitute. In low hydrogen export scenarios, onshore wind experiences high curtailment rates. However, at higher export volumes, increased system flexibility reduces curtailment rates.
    }
    \label{fig:el-cf}
\end{figure}


\begin{figure}
    \centering
        \begin{subfigure}[h]{0.49\textwidth}
        \centering
        \includegraphics[width=\textwidth]{graphics/heatmaps/cf-raw-ts-onshore wind.pdf}
    \end{subfigure}
    \begin{subfigure}[h]{0.49\textwidth}
        \centering
        \includegraphics[width=\textwidth]{graphics/heatmaps/cf-raw-ts-solar thermal.pdf}
    \end{subfigure}
    \begin{subfigure}[h]{0.49\textwidth}
        \centering
        \includegraphics[width=\textwidth]{graphics/heatmaps/cf-raw-ts-solar PV.pdf}
    \end{subfigure}
    \begin{subfigure}[h]{0.49\textwidth}
        \centering
        \includegraphics[width=\textwidth]{graphics/heatmaps/cf-raw-ts-run of river.pdf}
    \end{subfigure}
    \begin{subfigure}[h]{0.49\textwidth}
        \centering
        \includegraphics[width=\textwidth]{graphics/heatmaps/cop-ts-air-sourced heat pump.pdf}
    \end{subfigure}
    \begin{subfigure}[h]{0.49\textwidth}
        \centering
        \includegraphics[width=\textwidth]{graphics/heatmaps/cop-ts-ground-sourced heat pump.pdf}
    \end{subfigure}
    \caption{Capacity factors of renewable sources, and Coefficient of Performance (COP) of air-sourced and ground-sourced heat pumps.
    While solar thermal and solar photovoltaics (PV) show a diurnal pattern, onshore wind, run of river, and heat pumps have a seasonal pattern.
    }
    \label{fig:ren-cfs}
\end{figure}


\begin{figure*}[h] % Use 'figure*' to span both columns
    \centering
    \begin{subfigure}[b]{0.45\linewidth}
        \centering
        \includegraphics[width=\linewidth]{graphics/integrated_comp/contour_electrolysis_p_nom_opt_20_filterTrue_nFalse_exp200.pdf}
        \caption{Electrolysis optimal capacity}
        \label{fig:ely-p-nom-opt}
    \end{subfigure}
    \hfill
    \begin{subfigure}[b]{0.45\linewidth}
        \centering
        \includegraphics[width=\linewidth]{graphics/integrated_comp/contour_ft_p_nom_opt_20_filterTrue_nFalse_exp200.pdf}
        \caption{Fischer-Tropsch optimal capacity}
        \label{fig:ft-p-nom-opt}
    \end{subfigure}

    \hfill

    \caption{Electrolysis and Fischer-Tropsch capacities.
        Electrolysis capacity scales with hydrogen exports, but is also required for deep decarbonization scenarios above 80\% of \co mitigation. In those scenarios, Fischer-Tropsch capacities supply renewable fuels for the domestic demand.
    }
    \label{fig:ely-ft-p-nom-opt}
\end{figure*}

\clearpage